\chapter*{Abstract}
\addcontentsline{toc}{chapter}{Abstract}
LM-IDNet is a lightweight, unsupervised anomaly-detection system for identifying
abnormal network activity in IoT devices without attack signatures or
deep-learning models. The central idea is that each device produces
a characteristic mixture of network traffic—represented by TCP, UDP, SSDP,
and ARP packet counts—and that deviations from this pattern can provide
evidence of abnormal behaviour. LM-IDNet learns a normal device profile using
a Dirichlet–multinomial model and identifies statistically unlikely traffic
windows through calibrated likelihood scoring, complemented by compact volume
and ARP checks. The system integrates the Languasco–Migliardi numerical method
for accurate likelihood computation and supports guarded updates to the model
and decision thresholds. Experiments on smart-camera and Chromecast traffic
demonstrate high precision in the evaluated scenarios, while each alert retains
traceable statistical evidence. Resource measurements in a constrained
execution environment support the feasibility of lightweight operation. By
combining probabilistic modelling, accurate numerical computation, and
controlled adaptation in a reproducible pipeline, this thesis demonstrates an
interpretable approach to IoT intrusion detection that learns from normal
traffic without requiring labelled attack examples.
